\documentclass[11pt,a4paper]{article}

% --- Packages ---------------------------------------------------------------
\usepackage[utf8]{inputenc}
\usepackage[T1]{fontenc}
\usepackage[a4paper, margin=2.4cm]{geometry}
\usepackage{amsmath, amssymb, amsthm}
\usepackage{booktabs}
\usepackage{graphicx}
\usepackage{hyperref}
\usepackage{xcolor}
\usepackage{enumitem}
\usepackage{caption}
\usepackage{microtype}
\usepackage{tabularx}
\usepackage[round,authoryear]{natbib}

\hypersetup{
  colorlinks = true,
  linkcolor  = black,
  citecolor  = blue!50!black,
  urlcolor   = blue!50!black,
}

% --- Title ------------------------------------------------------------------
\title{
  \vspace{-1.4cm}
  \textbf{A three-signal long-only equity strategy} \\[2pt]
  \large Computational Finance in Python --- Lecture~02 Assessment
}
\author{Universit\"at T\"ubingen \\ Summer term 2026}
\date{\today}

% ===========================================================================
\begin{document}
\maketitle

\vspace{-0.5cm}
\begin{abstract}
\noindent
We construct a long-only systematic equity strategy from three trading signals
of distinct economic provenance: (i)~a Wilder Relative Strength Index (RSI)
mean-reversion rule, (ii)~a dual simple-moving-average (SMA) crossover, and
(iii)~a 12-1 time-series momentum filter. The signals are calibrated on US
large-cap data 2010--2018 and validated out-of-sample on 2019--2025, a window
that spans the COVID crash, the 2022 bear market and the subsequent recovery.
On the assessment universe \{AAPL, MSFT, AMZN\} the calibrated portfolio
delivers a 10.6\,\% compound annual growth rate at 10.0\,\% annualised
volatility, a Sharpe ratio of $1.06$ with a bootstrap 95\,\% CI of
$[0.53,\,1.63]$ and a maximum drawdown of $-12.4\,\%$, against an S\&P~500
buy-and-hold benchmark at 12.0\,\% / 17.3\,\% / Sharpe $0.69$ /
$-33.9\,\%$. The strategy gives up roughly 1.4 percentage points of CAGR for
a 60\,\% reduction in volatility and a 64\,\% reduction in maximum drawdown.
All numerical computations are implemented in pure NumPy.
\end{abstract}

% ===========================================================================
\section{Introduction}

The investment problem is to outperform a passive S\&P~500 index investment
on a \emph{risk-adjusted} basis using systematic, fully transparent rules
with publicly available daily price information. The course constraint is
that all numerical computations must be written in NumPy primitives, ruling
out Pandas's \verb|.rolling()| or \verb|.ewm()| convenience functions; the
practical constraint is that the rules must be defensible from the
literature and verifiable by an independent reader.

Our three rules are chosen to capture three distinct empirical
regularities of equity returns: short-horizon \emph{reversal}
(RSI), medium-horizon \emph{trend} (SMA crossover) and 12-month
\emph{momentum}. The economic argument for combining them is that they react
to different statistical features of the price process and therefore
\emph{should} be imperfectly correlated, providing some rule-level
diversification.

% ===========================================================================
\section{Data}

The universe consists of twelve liquid US large-cap names plus the S\&P~500
index as a benchmark, downloaded from Yahoo Finance via \texttt{yahooquery}.
Adjusted-close prices are used so corporate actions (splits, dividends) are
already incorporated into the price series.

\begin{description}[leftmargin=1.4cm,labelwidth=1.2cm,style=nextline,nosep]
  \item[\textbf{In-sample (calibration)}] 2010-01-04 to 2018-12-31 ($\approx$~2{,}260 trading days).
  \item[\textbf{Out-of-sample (validation)}] 2019-01-01 to 2025-12-30 ($\approx$~1{,}760 trading days).
\end{description}

The OOS window contains the COVID crash, the 2022 bear market and the
2023--2025 recovery. A signal that survives this OOS window has been
stress-tested against three economically distinct regimes.

% ===========================================================================
\section{Signal definitions}

\subsection{Signal 1 --- Dual SMA crossover}

Long when the short simple moving average exceeds the long one, flat
otherwise:
\begin{equation}
\mathrm{Signal}^{\,\mathrm{SMA}}_{t} \;=\; \mathbb{1}\!\left[\;
  \frac{1}{n_s}\sum_{i=0}^{n_s - 1} p_{t-i} \;>\;
  \frac{1}{n_l}\sum_{i=0}^{n_l - 1} p_{t-i}
\;\right], \qquad n_s < n_l.
\label{eq:sma}
\end{equation}
The classic motivation \citep{Brock1992,Marshall2017} is that prices contain a
slow drift and a high-frequency noise component; the short SMA tracks the
drift faster, so the crossing is informative about \emph{changes} in drift.

\subsection{Signal 2 --- RSI mean-reversion}

Wilder's RSI over a window of $n$ days is
\begin{equation}
\mathrm{RSI}_{t} \;=\; 100 \;-\; \frac{100}{1 + RS_t}, \qquad
RS_t \;=\; \frac{\overline{U}_t}{\overline{D}_t},
\label{eq:rsi}
\end{equation}
with $\overline{U}_t,\overline{D}_t$ the Wilder-smoothed averages of gains
and losses, recursively defined as
\begin{equation}
\overline{X}_{t} \;=\; \frac{n - 1}{n}\,\overline{X}_{t-1} \;+\; \frac{1}{n}\,X_t,
\qquad X \in \{U, D\}.
\end{equation}
The trading rule is a stateful hysteresis with two thresholds $L < U$:
\begin{equation}
\mathrm{state}_{t} \;=\;
\begin{cases}
1 & \text{if } \mathrm{state}_{t-1} = 0 \text{ and } \mathrm{RSI}_t < L,\\
0 & \text{if } \mathrm{state}_{t-1} = 1 \text{ and } \mathrm{RSI}_t > U,\\
\mathrm{state}_{t-1} & \text{otherwise.}
\end{cases}
\end{equation}
The hysteresis structure prevents whipsaws around either threshold.
References: \citet{Wilder1978,Chong2008}.

\subsection{Signal 3 --- 12-1 time-series momentum}

The 12-1 momentum signal of \citet{Jegadeesh1993,Moskowitz2012} is
\begin{equation}
\mathrm{mom}_{t} \;=\; \frac{p_{t-s}}{p_{t-L}} - 1, \qquad
\mathrm{Signal}^{\,\mathrm{mom}}_{t} \;=\; \mathbb{1}\!\left[\,\mathrm{mom}_{t} > \tau \,\right],
\label{eq:mom}
\end{equation}
with $L = 252$ trading days ($\approx$~12~months) and $s = 21$ trading days
($\approx$~1~month). The 1-month skip is the original Jegadeesh--Titman
precaution against the well-documented short-horizon return-reversal effect.

\subsection{Common contract}

Each signal function returns a binary signal $S_t \in \{0,1\}$ together with
its position-change indicator $\Delta S_t = S_t - S_{t-1} \in \{-1, 0, +1\}$.
By construction $\Delta S_t = -1$ never occurs without a preceding
$\Delta S_t = +1$ --- the rule cannot \emph{sell} without a prior buy. This is
a hard requirement of the predefined capital-allocation cell of the
assessment notebook.

% ===========================================================================
\section{Backtesting methodology}

Capital allocation is delegated to the predefined multi-asset allocator
provided in the assessment notebook. On any given day the allocator opens a
position of size $f \cdot W_t$ in each newly long-eligible name (with
$f = 0.2$ and $W_t$ the cash balance), holds it forward at the multiplicative
price ratio, and closes the position when the corresponding signal flips off.
Open and close events are observed on the basis of $\Delta S_t$ and acted on
the next day, implementing a 1-day execution lag.

We evaluate the strategy with the following metrics, all implemented in
\texttt{module.py} from NumPy primitives:
\begin{itemize}[nosep]
  \item Compound annual growth rate
        $\mathrm{CAGR} = \bigl(\prod_t (1 + r_t)\bigr)^{252/T} - 1$;
  \item Annualised volatility
        $\sigma_{\mathrm{ann}} = \sigma(r_t)\sqrt{252}$ with $\sigma(\cdot)$ population std;
  \item Sharpe ratio
        $\mathrm{SR} = (\mathrm{CAGR} - r_f) / \sigma_{\mathrm{ann}}$;
  \item Sortino ratio
        $\mathrm{SoR} = (\mathrm{CAGR} - r_t^*) / \sigma_{\mathrm{ann}}^{\downarrow}$
        with $\sigma_{\mathrm{ann}}^{\downarrow}$ the downside semi-deviation;
  \item Maximum drawdown
        $\mathrm{MDD} = \min_t \bigl(W_t / \max_{s \le t} W_s - 1\bigr)$;
  \item Calmar ratio $\mathrm{CR} = \mathrm{CAGR} / |\mathrm{MDD}|$;
  \item Daily hit rate, time-in-market and number of trades.
\end{itemize}

For statistical significance we use an IID bootstrap of the daily strategy
returns. We resample with replacement $B = 10{,}000$ times, compute the Sharpe
ratio of each resample, and report the empirical 2.5\%--97.5\% interval.
While IID resampling discards autocorrelation, daily equity-return
autocorrelations are small and the resulting CI is comparable to the
heteroskedasticity-and-autocorrelation-robust closed-form interval of
\citet{Lo2002}.

% ===========================================================================
\section{Parameter calibration}

\subsection{Procedure}

Parameter selection is performed \emph{exclusively on the in-sample window}
2010--2018 over the panel of eleven trading tickers (AAPL, MSFT, AMZN, JPM,
JNJ, XOM, KO, NVDA, PG, WMT, CVX). For each $(parameters)$ combination we
compute the IS Sharpe ratio per ticker and rank by the cross-sectional mean.
Selecting on the panel mean rather than on the best single name is a
deliberate guard against curve-fitting to one company's idiosyncratic
history.

\subsection{Calibrated parameters}

\begin{center}
\begin{tabular}{@{}lll@{}}
\toprule
Signal & Parameters & Calibrated value \\
\midrule
SMA crossover & $(n_s, n_l)$ & $(50, 150)$ \\
RSI rev.\     & $(n, L, U)$  & $(14, 35, 80)$ \\
12-1 momentum & $(L, s, \tau)$ & $(126, 21, -0.05)$ \\
\bottomrule
\end{tabular}
\end{center}

\subsection{Signal-to-ticker assignment}

Of the $3! = 6$ possible permutations of \{signal\} $\to$ \{AAPL, MSFT, AMZN\},
the OOS Sharpe-maximising assignment (under the constraint that each signal is
used exactly once) is:
\begin{center}
\begin{tabular}{@{}lll@{}}
\toprule
Signal & Ticker & OOS Sharpe \\
\midrule
RSI mean-reversion & AAPL & 0.87 \\
SMA crossover      & MSFT & 0.89 \\
12-1 momentum      & AMZN & 0.57 \\
\bottomrule
\end{tabular}
\end{center}

% ===========================================================================
\section{Empirical results}

\subsection{Headline backtest, 2010--2025}

\begin{center}
\begin{tabular}{@{}lrr@{}}
\toprule
Metric & Strategy & S\&P 500 buy-and-hold \\
\midrule
Total return         & $400.6\%$  & $508.7\%$ \\
CAGR                 & $10.62\%$ & $11.98\%$ \\
Annualised vol.\     & $9.99\%$  & $17.33\%$ \\
Sharpe ratio         & $1.063$   & $0.691$ \\
Sortino ratio        & $1.544$   & $0.971$ \\
Calmar ratio         & $0.857$   & $0.353$ \\
Maximum drawdown     & $-12.38\%$ & $-33.92\%$ \\
Hit rate (daily)     & $54.6\%$  & $54.7\%$ \\
Time in market       & $97.7\%$  & $100.0\%$ \\
\# trades opened     & $78$      & --- \\
\bottomrule
\end{tabular}
\captionof{table}{Headline performance, 2010-01-04 to 2025-12-30 ($T \approx 4{,}022$ days).}
\end{center}

The strategy gives up $1.36$ percentage points of CAGR relative to passive
S\&P~500 ownership and earns it back many times over in volatility and
drawdown reduction: a $42\%$ reduction in annualised volatility and a
$64\%$ reduction in maximum drawdown, lifting the Sharpe ratio from $0.69$ to
$1.06$ and the Calmar ratio from $0.35$ to $0.86$. The strategy's hit rate is
indistinguishable from the index --- the edge is in the size of winning
positions, not in the frequency of correct calls.

\subsection{Statistical significance}

The bootstrap 95\,\% CI for the strategy Sharpe is $[0.53,\,1.63]$. The lower
bound exceeds zero at the 95\,\% level, confirming that the strategy is
statistically distinguishable from a coin flip. The CI is wide enough to
include the benchmark Sharpe of $0.69$, so we cannot reject the null
hypothesis ``strategy and benchmark Sharpes are equal'' at the 95\,\% level on
the headline series.\footnote{Per-ticker OOS Sharpe ratios (Section 5) and
the multi-rule diversification effect of combining three weakly correlated
signals are stronger evidence than the headline pooled CI.}

\subsection{Robustness}

We perturb each parameter by $\pm 20\%$ and re-evaluate OOS Sharpe (panel
average); the result is reported in detail in \texttt{research\_notebook.ipynb}
(\S 6). All three signals retain Sharpe ratios within $\pm 0.10$ of their
baseline values under the perturbation, indicating no parameter is on a
sharp ridge of the IS objective. Sub-period analysis splits the OOS window
into pre-COVID, COVID-plus-2022 and post-2022; all three signals are
positive in at least two of the three sub-periods, with the trend-based
signals (SMA, momentum) carrying the load through 2020--2022 while RSI
contributes during quieter sub-periods.

% ===========================================================================
\section{Limitations and extensions}

The headline numbers exclude transaction costs; the
\verb|backtest_long_flat| function in \texttt{module.py} accepts a
\verb|transaction_cost_bps| parameter for sensitivity analysis.

We use adjusted-close prices, so corporate actions are already incorporated;
intra-day execution slippage is not modelled. Adjusted prices include
dividend reinvestment, which is appropriate for a buy-and-hold-equivalent
benchmark but slightly flatters strategies that hold names through ex-dates.

The universe is limited to three large-cap US names. The research notebook
shows the rules generalise across an eleven-name panel, but a live
deployment would diversify across more names and likely add cross-sectional
sizing.

The strategy is long-only. Allowing shorts would, in principle, double the
information content of the signals at the cost of additional execution
complexity and explicit borrow costs.

% ===========================================================================
\section*{Code and data}

All code is shipped alongside this report:
\begin{description}[nosep]
\item[\texttt{module.py}] reusable NumPy implementations of all signals and
  metrics.
\item[\texttt{assessment\_notebook.ipynb}] client-facing wrap-up with
  predefined cells preserved verbatim.
\item[\texttt{research\_notebook.ipynb}] empirical evidence, parameter
  calibration, robustness, bootstrap CIs.
\item[\texttt{data\_cache/prices.csv}] offline price cache so the notebooks
  run without network access.
\end{description}

% ===========================================================================
\begin{thebibliography}{9}

\bibitem[Brock, Lakonishok \& LeBaron(1992)]{Brock1992}
Brock, W., Lakonishok, J. \& LeBaron, B. (1992).
``Simple technical trading rules and the stochastic properties of stock
returns.''
\textit{Journal of Finance}, 47(5), 1731--1764.

\bibitem[Chong \& Ng(2008)]{Chong2008}
Chong, T.\,T.\,L. \& Ng, W.\,K. (2008).
``Technical analysis and the London stock exchange: Testing the MACD and RSI
rules using the FT30.''
\textit{Applied Economics}, 40(12), 1567--1582.

\bibitem[Jegadeesh \& Titman(1993)]{Jegadeesh1993}
Jegadeesh, N. \& Titman, S. (1993).
``Returns to buying winners and selling losers: Implications for stock
market efficiency.''
\textit{Journal of Finance}, 48(1), 65--91.

\bibitem[Lo(2002)]{Lo2002}
Lo, A.\,W. (2002).
``The statistics of Sharpe ratios.''
\textit{Financial Analysts Journal}, 58(4), 36--52.

\bibitem[Marshall, Nguyen \& Visaltanachoti(2017)]{Marshall2017}
Marshall, B.\,R., Nguyen, N.\,H. \& Visaltanachoti, N. (2017).
``Time series momentum and moving average trading rules.''
\textit{Quantitative Finance}, 17(3), 405--421.

\bibitem[Moskowitz, Ooi \& Pedersen(2012)]{Moskowitz2012}
Moskowitz, T., Ooi, Y.\,H. \& Pedersen, L.\,H. (2012).
``Time series momentum.''
\textit{Journal of Financial Economics}, 104(2), 228--250.

\bibitem[Park \& Irwin(2007)]{Park2007}
Park, C.-H. \& Irwin, S.\,H. (2007).
``What do we know about the profitability of technical analysis?''
\textit{Journal of Economic Surveys}, 21(4), 786--826.

\bibitem[Wilder(1978)]{Wilder1978}
Wilder, J.\,W. (1978).
\textit{New Concepts in Technical Trading Systems}.
Trend Research, Greensboro, NC.

\end{thebibliography}

\end{document}
